\appendix
\chapter{Lean Setup and Conventions}
\label{app:lean-setup}

\section{Minimal Project Setup}

Every listing in this book is ordinary Lean~4, checked against
\texttt{mathlib} for the handful of lemmas the later chapters borrow
(\texttt{omega}, \texttt{Set.mem\_setOf\_eq}). A reader following
along needs no more than the following two files in a fresh
directory managed by \texttt{lake}, Lean's build tool.

\texttt{lean-toolchain} pins the compiler version so that a project
built today keeps working:
\begin{lean}
leanprover/lean4:v4.11.0
\end{lean}

\texttt{lakefile.toml} declares the project and its one dependency:
\begin{lean}
name = "proof_and_countermodel"
defaultTargets = ["ProofAndCountermodel"]

[[require]]
name = "mathlib"
scope = "leanprover-community"

[[lean_lib]]
name = "ProofAndCountermodel"
\end{lean}
Running \texttt{lake update} once, followed by \texttt{lake build},
fetches \texttt{mathlib} and compiles it; this is the slow step, and
only needs to happen once per toolchain version. Every
\texttt{example}/\texttt{theorem} block quoted in this book can then
be pasted into a single \texttt{.lean} file under
\texttt{ProofAndCountermodel/} and checked with
\texttt{lake env lean ThatFile.lean}, or browsed interactively with
any Lean-aware editor (VS Code with the \texttt{lean4} extension is
the reference setup used while writing this book).

\section{\texttt{Prop} versus \texttt{Bool}}

Chapter~1 introduces two different Lean citizens that both claim to
speak for the five familiar connectives, and the distinction is worth
restating plainly here since every later chapter leans on it without
re-explaining it.

A term of type \texttt{Prop} is a \emph{proposition}: a statement
that may or may not be true, whose type is inhabited exactly when it
is provable. A proof of $p \land q$ is not a fact looked up
somewhere — it is a concrete piece of data, an \texttt{And.intro}
structure holding a proof of $p$ and a proof of $q$ side by side.
Two different proofs of the same proposition need not be equal as
terms; Lean does not ask them to be, since \texttt{Prop} is
proof-irrelevant only up to definitional equality, and the book never
needs anything stronger than that.

A term of type \texttt{Bool}, by contrast, is a \emph{computed
value}: exactly \texttt{true} or exactly \texttt{false}, produced by
running a decidable total function. \texttt{decide} is the tactic
that crosses from one register to the other: given a decidable
\texttt{Prop}, it evaluates the corresponding \texttt{Bool}
computation and, if that computation reduces to \texttt{true}, uses
the reduction itself as the proof. This is why every
\texttt{by decide} in this book is doing two things at once — settling
a finite question by brute enumeration, and silently converting that
enumeration into a \texttt{Prop}-level proof term.

The practical rule of thumb used throughout: reach for \texttt{Bool}
and \texttt{decide} whenever the domain is finite and the question is
a concrete instance (a specific \texttt{Fin n}, a specific truth
table); reach for \texttt{Prop} and an explicit proof term whenever
the claim is about \emph{all} propositions, all naturals, or any
other domain too large to enumerate. Chapters~1, 4, and 5 mix both
registers side by side precisely to make that boundary visible rather
than leaving it implicit.

\begin{thebibliography}{9}

\bibitem{avigad2021}
J.~Avigad, L.~de~Moura, S.~Kong, and S.~Ebner,
\emph{Theorem Proving in Lean 4},
online textbook, 2021.
Available online: \url{https://leanprover.github.io/theorem_proving_in_lean4/}.

\bibitem{vandalen2013}
D.~van Dalen,
\emph{Logic and Structure}, 5th ed.,
Universitext, Springer, 2013.

\bibitem{priest2008}
G.~Priest,
\emph{An Introduction to Non-Classical Logic: From If to Is}, 2nd ed.,
Cambridge University Press, 2008.

\end{thebibliography}
